% Part: computability
% Chapter: recursive-functions
% Section: pr-functions-computable

\documentclass[../../../include/open-logic-section]{subfiles}

\begin{document}

% SEGMENT OLP-0215-S01

\olfileid{cmp}{rec}{cmp}
\olsection{முதன்மை மீள்வரையறைச் சார்புகள் கணிக்கத்தக்கவை}

சார்பு $h$ பின்வரும் முதன்மை மீள்வரையறையால் வரையறுக்கப்படுகிறது என்று
கொள்க:
\begin{eqnarray*}
h(\vec x, 0)     & = & f(\vec x) \\
h(\vec x, y + 1) & = & g(\vec x, y, h(\vec x, y))
\end{eqnarray*}
மேலும் சார்புகள் $f$ மற்றும் $g$ கணிக்கத்தக்கவை என்றும் கொள்க.
($x_0$, \dots, $x_{k-1}$ ஆகியவற்றைச் சுருக்கமாகக் குறிக்க
$\vec x$ ஐப் பயன்படுத்துகிறோம்.) அப்போது $h(\vec x, 0)$ ஐத்
தெளிவாகவே கணிக்கலாம்; ஏனெனில் அது $f(\vec x)$ மட்டுமே, அதைக்
கணிக்கத்தக்கது என்று முன்கொண்டுள்ளோம். $h(\vec x, 1)$ ஐயும்
கணிக்கலாம்; ஏனெனில் $1 = 0 + 1$; எனவே $h(\vec x, 1)$ என்பது
பின்வருவதே:
\begin{align*}
 h(\vec x, 1) & = g(\vec x, 0, h(\vec x, 0)) =  g(\vec x, 0, f(\vec x)).
\intertext{இம்முறையில் தொடர்ந்து பின்வருவனவற்றைக் கணிக்கலாம்:}
h(\vec x, 2) & = g(\vec x, 1, h(\vec x, 1)) = g(\vec x, 1, g(\vec x, 0, f(\vec x)))\\
h(\vec x, 3) & = g(\vec x, 2, h(\vec x, 2)) = g(\vec x, 2, g(\vec x, 1, g(\vec x, 0, f(\vec x))))\\
h(\vec x, 4) & = g(\vec x, 3, h(\vec x, 3)) = g(\vec x, 3, g(\vec x, 2, g(\vec x, 1, g(\vec x, 0, f(\vec x)))))\\
& \vdots
\end{align*}
எனவே, பொதுவாக $h(\vec x, y)$ ஐக் கணிக்க, முதலில்
$h(\vec x, 0)$, பின்னர் $h(\vec x, 1)$, \dots, என அடுத்தடுத்து
$h(\vec x, y)$ ஐ அடையும் வரை கணிக்க.

ஆகவே சார்புகள் $f$ மற்றும் $g$ கணிக்கத்தக்கவை எனில், முதன்மை
மீள்வரையறை ஒரு புதிய கணிக்கத்தக்க சார்பை விளைவிக்கிறது. சார்புகள்
$f$ மற்றும் $g_i$ கணிக்கத்தக்கவை எனில், சார்புகளின் சேர்ப்பும்
ஒரு கணிக்கத்தக்க சார்பை விளைவிக்கிறது.

அடிப்படைச் சார்புகள் $\Zero$, $\Succ$, $\Proj{n}{i}$ ஆகியவை
கணிக்கத்தக்கவை; மேலும் சேர்ப்பும் முதன்மை மீள்வரையறையும் கணிக்கத்தக்க
சார்புகளிலிருந்து கணிக்கத்தக்க சார்புகளை விளைவிக்கின்றன. எனவே ஒவ்வொரு
முதன்மை மீள்வரையறைச் சார்பும் கணிக்கத்தக்கது.

\end{document}
